\documentclass[12pt,a4paper]{article}

\usepackage[utf8]{inputenc}
\usepackage[indonesian]{babel}
\usepackage{amsmath,amssymb,amsfonts}
\usepackage{geometry}
\usepackage{booktabs}
\usepackage{array}
\usepackage{hyperref}

\geometry{
    top=2.5cm,
    bottom=2.5cm,
    left=3cm,
    right=2.5cm
}

\title{\textbf{Slope dan Offset pada Kalibrasi Sensor pH:\\Teori, Perhitungan, dan Penggunaan di Firmware}}
\author{Tim Riset Glicovia}
\date{20 Oktober 2026}

\begin{document}

\maketitle

\section{Persamaan Nernst}

Elektroda pH menghasilkan tegangan $E$ yang bergantung pada konsentrasi ion $\text{H}^+$:
\begin{equation}
E = E^0 - \frac{2.303 \cdot R \cdot T}{F} \cdot \text{pH}
\end{equation}

Persamaan ini berbentuk garis lurus $y = m \cdot x + c$, dengan pemetaan pada Tabel~\ref{tab:simbol}.

\begin{table}[h]
\centering
\begin{tabular}{@{}llp{6.5cm}@{}}
\toprule
Simbol & Di sensor & Arti \\
\midrule
$y$ & tegangan $V$ yang dibaca ADC & keluaran sensor \\
$x$ & pH & besaran yang ingin dicari \\
$m$ (\texttt{PH\_SLOPE}) & $-2.303\,R\,T/F$ & besar perubahan tegangan per 1 unit pH \\
$c$ (\texttt{PH\_OFFSET}) & $E^0$ & tegangan saat pH $= 0$ (titik potong sumbu $V$) \\
\bottomrule
\end{tabular}
\caption{Pemetaan persamaan Nernst ke persamaan garis lurus.}
\label{tab:simbol}
\end{table}

\section{Arti Slope dan Offset}

\paragraph{Slope ($m$)} adalah kemiringan garis. Secara teori, pada $25^\circ\text{C}$:
\begin{equation}
m = \frac{2.303 \times 8.314 \times 298.15}{96485} = 0.05916\ \text{V/pH} = 59.16\ \text{mV/pH}
\end{equation}
Tandanya negatif karena makin asam (pH kecil), tegangan makin tinggi. Pada SPCE, slope nyata biasanya lebih kecil (\textit{sub-Nernstian}), misalnya $-45$ sampai $-58\ \text{mV/pH}$. Karena itu sensor harus dikalibrasi, bukan memakai nilai $59.16\ \text{mV/pH}$ secara langsung.

\paragraph{Offset ($c$)} menggeser garis naik atau turun. Nilainya bergantung pada bahan elektroda, elektroda referensi, dan rangkaian penguat, sehingga berbeda untuk setiap sensor.

\section{Cara Menghitung}

Ukur larutan penyangga (\textit{buffer}) yang pH-nya diketahui, lalu cari garis yang melewati titik-titik tersebut. Landasannya adalah geometri analitik (garis lurus) dan regresi linier kuadrat terkecil.

\subsection{Dua titik}

Dari titik $(\text{pH}_1, V_1)$ dan $(\text{pH}_2, V_2)$:
\begin{align}
m &= \frac{V_2 - V_1}{\text{pH}_2 - \text{pH}_1} \label{eq:m2}\\
c &= V_1 - m \cdot \text{pH}_1 \label{eq:c2}
\end{align}

\paragraph{Contoh angka (hipotetis, bukan data pengukuran).}
Misalkan pH $4.01$ terbaca $0.2000\ \text{V}$ dan pH $7.00$ terbaca $0.0300\ \text{V}$. Maka
\begin{align*}
m &= \frac{0.0300 - 0.2000}{7.00 - 4.01} = \frac{-0.1700}{2.99} = -0.05686\ \text{V/pH} \quad (-56.9\ \text{mV/pH})\\
c &= 0.2000 - (-0.05686 \times 4.01) = 0.2000 + 0.2280 = 0.4280\ \text{V}
\end{align*}
Pemeriksaan untuk pH $7.00$: $-0.05686 \times 7 + 0.4280 = 0.0300\ \text{V}$.

\subsection{Tiga titik atau lebih (regresi)}

Dicari $m$ dan $c$ yang meminimalkan jumlah kuadrat selisih
\begin{equation}
\sum_i \bigl(V_i - (m\,\text{pH}_i + c)\bigr)^2
\end{equation}
dengan solusi
\begin{align}
m &= \frac{\sum_i (\text{pH}_i - \overline{\text{pH}})(V_i - \overline{V})}{\sum_i (\text{pH}_i - \overline{\text{pH}})^2} \\
c &= \overline{V} - m\,\overline{\text{pH}}
\end{align}
Perhitungan ini dilakukan oleh \texttt{np.polyfit(ph, volt, 1)} pada \texttt{calibrate.py}. Titik ketiga berguna untuk memeriksa kelinearan sensor melalui koefisien determinasi $R^2$.

\section{Penggunaan di Firmware}

Pembacaan $V$ dari kanal A2 ADS1115 diubah menjadi pH dengan membalik persamaan garis:
\begin{equation}
\text{pH} = \frac{V - c}{m}
\end{equation}
Pada contoh di atas, sampel yang terbaca $0.1000\ \text{V}$ menghasilkan
\[
\text{pH} = \frac{0.1000 - 0.4280}{-0.05686} = 5.77
\]

\section{Catatan Teori}

\begin{itemize}
    \item Persamaan Nernst mengasumsikan suhu $25^\circ\text{C}$. Slope berubah sekitar $0.2\ \text{mV/pH}$ per $^\circ\text{C}$. Jika suhu saliva ($\approx 35^\circ\text{C}$) berbeda jauh dari suhu kalibrasi, timbul galat kecil. Untuk akurasi lebih baik, kalibrasi dilakukan pada suhu yang mendekati kondisi pemakaian.
    \item Garis lurus hanya berlaku di antara titik kalibrasi. Di luar rentang itu, hasil merupakan ekstrapolasi.
    \item Jika slope hasil kalibrasi jauh dari $-59\ \text{mV/pH}$ (misalnya $-5\ \text{mV/pH}$), sensor tidak merespons pH dengan baik.
\end{itemize}

\end{document}
